% !PRINTED 276
\weberpage{276}
Es folgt zunächst aus \S 69, (9):
\[
c_1 \equiv c+\partial \pmod{a}, \quad c_1' \equiv c'+\partial' \pmod{a'}, \quad c_1'' \equiv c''+\partial'' \pmod{a''},
\]
woraus, nach (4)
\[
\begin{aligned}
\partial''c_1' &\equiv \partial''c' + \partial \equiv c_1 \pmod{a'}\\
\partial'c_1'' &\equiv \partial'c'' + \partial \equiv c_1 \pmod{a''}
\end{aligned}
\]
in Übereinstimmung mit den Kongruenzen (4).

Für die Zusammensetzung (6) ergibt sich nach \S 69, (12), daß
$\partial_2$, $\partial_2'$, $\partial_2''$ die größten gemeinschaftlichen Teiler von $a,c$;
$a',c'$; $a'',c''$ sind; weil aber $\partial''$ relativ prim zu $a'$, und
$\partial''c' \equiv c \pmod{a'}$ ist, so ist auch $\partial_2'$ der größte gemeinschaftliche
Teiler von $a'$ und $c$ und aus den gleichen Gründen $\partial_2''$ der größte
gemeinschaftliche Teiler von $a''$ und $c$, woraus man schließt, da
$a'$, $a''$ relativ prim sind:
\[
\partial_2 = \partial_2'\partial_2'', \qquad a_2 = a_2'a_2''.
\]
Ferner ist nach \S 69, (11), (12), (13):
\[
\begin{aligned}
\partial_2 &= \alpha c - \gamma a, & c_2 &= -\partial\alpha,\\
\partial_2' &= \alpha'c' - \gamma'a', & c_2' &= -\partial'\alpha',
\end{aligned}
\]
also nach (4):
\[
\partial_2c_2' - c_2\partial_2' \equiv \alpha\alpha'(\partial c' - c\partial') \equiv 0 \pmod{n'},
\]
folglich auch
\[
\partial_2''c_2' - c_2 \equiv 0 \pmod{a_2'},
\]
in Übereinstimmung mit (4), und ebenso folgt:
\[
\partial_2'c_2'' - c_2 \equiv 0 \pmod{a_2''},
\]
wodurch also der Beweis geführt ist, daß die durch (4) ausgedrückte
Zusammengehörigkeit der Transformationen
\[
\begin{pmatrix} a, & 0\\ c, & \partial \end{pmatrix}, \quad
\begin{pmatrix} a', & 0\\ c', & \partial' \end{pmatrix}, \quad
\begin{pmatrix} a'', & 0\\ c'', & \partial'' \end{pmatrix}
\]
durch Anwendung irgend einer \gesperrt{linearen} Transformation auf $\omega$
nicht gestört wird. Da wir hier $n$ als ungerade voraussetzen, so
können wir immer $c$, $c'$, $c''$ durch 16 teilbar annehmen; und wenn $n$
und folglich auch $n'$, $n''$ durch 3 unteilbar sind, so können $c$, $c'$, $c''$
auch durch 3 teilbar angenommen werden. Ist aber $n$ durch 3
teilbar, so wird von den beiden Faktoren $n'$, $n''$ der eine, etwa $n''$,
durch 3 teilbar sein, der andere, $n'$, nicht. Es kann dann $c'$
noch durch 3 teilbar vorausgesetzt werden, nicht aber $c$ und $c''$.
In diesem Falle soll die Abhängigkeit des $c''$ von $c$ noch näher
bestimmt werden durch die Kongruenz
\weq{7}{
\partial'c'' \equiv c \pmod{3a''}.
}
